\sisection{S5. Calibration recovery and exploratory subspace compression}

S5.1 and S5.2 come from \path{analysis/calibration_panel_recovery.py}, with frozen outputs in
\path{results/calibration_panel_recovery/}. The matched random comparison in S5.2 comes from
\path{analysis/public_scaffold_holdout_recovery.py} and
\path{results/public_scaffold_holdout_recovery/}. S5.3 comes from
\path{analysis/biological_core_modes.py} and \path{results/biological_core_modes/}.

\sisubsection{S5.1 Recovery of the source-library map}

One replicate draws $n$ ligands without replacement from the complete row list, two-way-centres
the sampled score block, and takes the target correlation matrix of the centred block. Recovery
is the Spearman correlation between the upper triangle of that map and the upper triangle of the
full-matrix map: 946 target pairs for Docking-44 and 1{,}653 for DOCKSTRING-58. The draw reads no
experimental value and no score threshold. We ran 100 replicates at each of seven sizes, with
seed 20260803 for DOCKSTRING-58 and seed 20260804 for Docking-44, both recorded in
\path{results/calibration_panel_recovery/metadata.json}.

The estimand declared in that file is recovery of the full-matrix map, so a calibration sample is
a subset of the rows its own reference is built from. A separate bundle runs both reference
scopes side by side on the same design. Scoring each sample against the map of the rows it did
not use gives 0.940 at 200 ligands and 0.973 at 500 for Docking-44, against 0.941 and 0.975 when
the calibration rows stay in the reference; for DOCKSTRING-58 the four values are 0.920, 0.965,
0.920 and 0.966
(\path{results/revision_map_decision_sensitivities/recovery_edge_summary.csv}). The choice of
reference scope moves the mean by at most 0.002.

\begin{table}[htbp]
\centering
\caption{Recovery of the full-matrix residual target map from an outcome-blind calibration
sample. Entries are means over 100 replicates per size with 2.5--97.5\% repeated-sample ranges.
Source: \texttt{results/calibration\_panel\_recovery/summary.csv}. The reference residual
participation ratios are 9.303 for Docking-44 and 18.206 for DOCKSTRING-58.}
\label{tab:s5curve}
\small
\begin{tabular}{lrlrr}
\toprule
Matrix & $n$ & Map $\rho$ & Plug-in PR & OAS PR \\
\midrule
Docking-44 & 50 & 0.813 [0.657, 0.901] & 7.13 & 10.50 \\
Docking-44 & 100 & 0.898 [0.772, 0.946] & 7.83 & 9.70 \\
Docking-44 & 200 & 0.942 [0.898, 0.962] & 8.63 & 9.67 \\
Docking-44 & 500 & 0.976 [0.961, 0.984] & 9.03 & 9.47 \\
Docking-44 & 1{,}000 & 0.987 [0.977, 0.991] & 9.21 & 9.43 \\
Docking-44 & 2{,}000 & 0.994 [0.991, 0.996] & 9.24 & 9.35 \\
Docking-44 & 5{,}000 & 0.998 [0.997, 0.999] & 9.27 & 9.32 \\
\midrule
DOCKSTRING-58 & 50 & 0.766 [0.665, 0.822] & 12.80 & 18.70 \\
DOCKSTRING-58 & 100 & 0.861 [0.808, 0.896] & 14.94 & 18.10 \\
DOCKSTRING-58 & 200 & 0.918 [0.885, 0.938] & 16.48 & 18.15 \\
DOCKSTRING-58 & 500 & 0.965 [0.951, 0.972] & 17.42 & 18.11 \\
DOCKSTRING-58 & 1{,}000 & 0.982 [0.977, 0.986] & 17.90 & 18.25 \\
DOCKSTRING-58 & 2{,}000 & 0.991 [0.987, 0.992] & 18.01 & 18.19 \\
DOCKSTRING-58 & 5{,}000 & 0.996 [0.995, 0.997] & 18.15 & 18.22 \\
\bottomrule
\end{tabular}
\end{table}

Table~\ref{tab:s5curve} gives the full curve. Going from 50 to 200 ligands adds 0.129 for
Docking-44 and 0.152 for DOCKSTRING-58; going from 200 to 500 adds 0.034 and 0.047. A further
tenfold increase, from 500 to 5{,}000, adds 0.022 and 0.032. The replicate spread narrows faster
than the mean rises. At 50 ligands the 2.5--97.5\% range spans 0.244 for Docking-44 and 0.158 for
DOCKSTRING-58; at 200 ligands it spans 0.065 and 0.052.

The plug-in participation ratio of a small sample is biased downward, because finite-sample
correlation noise inflates the mean squared off-diagonal correlation. Oracle approximating
shrinkage removes most of that bias \cite{chen2010}. At 200 ligands the plug-in estimate is 8.63
against the Docking-44 reference of 9.303, while the shrunk estimate is 9.67. For DOCKSTRING-58
the two estimates are 16.48 and 18.15 against a reference of 18.206. Recovering the ordering of
the map and recovering its participation ratio are separate questions at these sizes.

\begin{table}[htbp]
\centering
\caption{Score cells needed by a two-stage calibration design. Full-matrix cells are ligands
times targets. Source: the \texttt{cost\_examples} block of
\texttt{results/calibration\_panel\_recovery/metadata.json}.}
\label{tab:s5cells}
\small
\begin{tabular}{lrrrlr}
\toprule
Matrix & Ligands & Targets & $n$ & Cells scored & Reduction \\
\midrule
Docking-44 & 12{,}651 & 44 & 200 & 8{,}800 of 556{,}644 (1.58\%) & 63.3-fold \\
Docking-44 & 12{,}651 & 44 & 500 & 22{,}000 of 556{,}644 (3.95\%) & 25.3-fold \\
DOCKSTRING-58 & 260{,}060 & 58 & 200 & 11{,}600 of 15{,}083{,}480 (0.077\%) & 1{,}300.3-fold \\
DOCKSTRING-58 & 260{,}060 & 58 & 500 & 29{,}000 of 15{,}083{,}480 (0.192\%) & 520.1-fold \\
\bottomrule
\end{tabular}
\end{table}

Table~\ref{tab:s5cells} gives the arithmetic for both matrices. Every calibration ligand is
scored against every target, so the reduction in cells equals the reduction in ligands and the
target count cancels. The two matrices differ by a factor of 20.6 in that reduction because
DOCKSTRING-58 has 20.6 times more rows.

Two downstream quantities were tracked on the same replicates. The first is retrieval of the
strongest edges of the full-matrix map: labelling the top decile of reference edges positive and
ranking by the sample's edge values gives, at 200 ligands, ROC-AUC 0.980 and average precision
0.862 for Docking-44, and 0.976 and 0.855 for DOCKSTRING-58. At 500 ligands the four values are
0.991, 0.930, 0.990 and 0.933. The second is the frozen 20-kinase analysis endpoint, in which
a target pair is positive when its centred experimental correlation falls in the upper decile of
at least two assay panels, labelling 15 of 190 pairs. That endpoint applies to DOCKSTRING-58
only, because the 20 kinases are DOCKSTRING columns.

\begin{table}[htbp]
\centering
\caption{The frozen 20-kinase experimental-pair endpoint, 15 positives of 190, evaluated from the
DOCKSTRING-58 calibration samples. Entries are means over 100 replicates with 2.5--97.5\% ranges.
The last row is the value from all 260{,}060 complete DOCKSTRING rows. Sources:
\texttt{results/calibration\_panel\_recovery/summary.csv} and the
\texttt{reference\_experimental\_pair\_metrics} block of the same directory's
\texttt{metadata.json}.}
\label{tab:s5endpoint}
\small
\begin{tabular}{lll}
\toprule
$n$ & ROC-AUC & Average precision \\
\midrule
50 & 0.677 [0.594, 0.757] & 0.229 [0.140, 0.350] \\
100 & 0.691 [0.595, 0.768] & 0.259 [0.180, 0.353] \\
200 & 0.701 [0.646, 0.772] & 0.269 [0.188, 0.344] \\
500 & 0.706 [0.672, 0.742] & 0.276 [0.213, 0.340] \\
1{,}000 & 0.708 [0.680, 0.743] & 0.293 [0.232, 0.341] \\
2{,}000 & 0.710 [0.686, 0.732] & 0.291 [0.242, 0.328] \\
5{,}000 & 0.710 [0.696, 0.722] & 0.296 [0.255, 0.318] \\
\midrule
260{,}060 & 0.710 & 0.301 \\
\bottomrule
\end{tabular}
\end{table}

Table~\ref{tab:s5endpoint} shows that this endpoint saturates earlier than the map itself.
ROC-AUC reaches 0.701 at 200 ligands against the full-matrix 0.710, and from 500 ligands upward
the replicate mean stays within 0.005 of that value. Average precision is noisier and not
monotone: the mean at 2{,}000 ligands, 0.291, sits below the mean at 1{,}000, 0.293. The
full-matrix average precision of 0.301 lies inside the 2.5--97.5\% range at every size from 200
upward, so these replicates cannot separate a 200-ligand sample from the complete matrix on that
metric.

These are score-surface diagnostics. A calibration sample recovers the map of the library it was
drawn from, and it neither validates a pose nor estimates an affinity.

\sisubsection{S5.2 Chemical-group-held-out recovery}

The holdout design splits each library into five \texttt{GroupKFold} folds of whole chemical
groups. Groups are the frozen \texttt{Butina\_clusters} labels for Docking-44, 8{,}150 clusters
over 12{,}651 rows \cite{butina1999}, and RDKit Bemis--Murcko scaffolds for DOCKSTRING-58 on a
fixed simple-random 15{,}000-row support drawn with seed 71, giving 11{,}905 groups
\cite{murcko1996,rdkit2025}. Calibration samples come only from the four training folds and the
reference map is estimated on the held-out fold, so no chemical group appears on both sides. The
producing script asserts an empty intersection of training and held-out group labels in every
split and raises on any overlap. We drew 25 samples per fold at 200 ligands and 25 at 500, giving
125 samples per matrix and size, with seed 20260805 for DOCKSTRING-58 and seed 20260806 for
Docking-44.

Per fold, Docking-44 holds out 2{,}530 or 2{,}531 ligands over 1{,}629 to 1{,}631 clusters, with
6{,}519 to 6{,}521 clusters left for calibration. DOCKSTRING-58 holds out 3{,}000 ligands over
2{,}357 to 2{,}387 scaffolds, with 9{,}518 to 9{,}548 left for calibration
(\path{results/calibration_panel_recovery/scaffold_holdout_metrics.csv}). The residual
participation ratio of the held-out fold itself moves from 8.72 to 10.23 across the Docking-44
folds and from 17.60 to 18.18 across the DOCKSTRING-58 folds, so the target of recovery is not
identical from fold to fold.

\begin{table}[htbp]
\centering
\caption{Recovery of a held-out chemical-group fold's residual map from calibration ligands in
the other four folds. Entries are the mean over 125 samples, the standard deviation, and the
observed extremes. Source:
\texttt{results/calibration\_panel\_recovery/scaffold\_holdout\_summary.csv}.}
\label{tab:s5group}
\small
\begin{tabular}{lrrrr}
\toprule
Matrix & $n$ & Mean $\rho$ & SD & Minimum--maximum \\
\midrule
Docking-44 & 200 & 0.927 & 0.023 & 0.830--0.964 \\
Docking-44 & 500 & 0.959 & 0.011 & 0.918--0.977 \\
DOCKSTRING-58 & 200 & 0.916 & 0.011 & 0.880--0.938 \\
DOCKSTRING-58 & 500 & 0.958 & 0.006 & 0.934--0.970 \\
\bottomrule
\end{tabular}
\end{table}

Table~\ref{tab:s5group} gives the result. Against the simple-random values of
Table~\ref{tab:s5curve}, holding out whole chemical groups costs 0.015 of mean agreement for
Docking-44 at 200 ligands and 0.002 for DOCKSTRING-58; at 500 ligands the two costs are 0.016 and
0.006. The worst single Docking-44 replicate at 200 ligands reaches 0.830, so the spread across
folds is wider than the gap between the two designs.

A separate bundle repeats the holdout against a matched random comparator that uses the same
calibration and evaluation sizes in each replicate but samples rows rather than groups. That
bundle applies a stricter grouping rule, canonical non-isomeric Bemis--Murcko scaffolds for
cyclic molecules and Standard-InChIKey connectivity blocks for acyclic molecules, and it
estimates Docking-44 target means fold-locally so that held-out scores never enter the
imputation. Its group-held-out means differ from Table~\ref{tab:s5group} by at most 0.003 in all
four rows.

\begin{table}[htbp]
\centering
\caption{Chemical-group holdout against a size-matched random-row holdout, 125 paired replicates
per row. Entries are means with 2.5--97.5\% repeated-sample ranges. Sources:
\texttt{scaffold\_holdout\_summary.csv},
\texttt{matched\_random\_row\_holdout\_summary.csv} and
\texttt{chemical\_group\_minus\_random\_summary.csv} in
\texttt{results/public\_scaffold\_holdout\_recovery/}.}
\label{tab:s5matched}
\small
\begin{tabular}{lrlll}
\toprule
Matrix & $n$ & Group-held-out $\rho$ & Matched-random $\rho$ & Paired difference \\
\midrule
Docking-44 & 200 & 0.925 [0.884, 0.956] & 0.935 [0.884, 0.960] & $-0.010$ [$-0.057$, 0.041] \\
Docking-44 & 500 & 0.958 [0.926, 0.975] & 0.968 [0.943, 0.980] & $-0.011$ [$-0.050$, 0.019] \\
DOCKSTRING-58 & 200 & 0.914 [0.889, 0.933] & 0.917 [0.893, 0.934] & $-0.003$ [$-0.034$, 0.029] \\
DOCKSTRING-58 & 500 & 0.958 [0.945, 0.965] & 0.960 [0.950, 0.968] & $-0.002$ [$-0.018$, 0.011] \\
\bottomrule
\end{tabular}
\end{table}

Every paired central range in Table~\ref{tab:s5matched} includes zero. All four point estimates
are negative, so the group split is the harder of the two by 0.010 and 0.011 for Docking-44 and
by 0.003 and 0.002 for DOCKSTRING-58. These are composition ranges over repeated splits on one
fixed target panel. They are not equivalence tests, and we do not read them as evidence that the
two designs are interchangeable.

The group split is harder by a structural diagnostic as well. For every held-out ligand we
computed the exact maximum Tanimoto similarity to the whole calibration pool on achiral
radius-2, 2{,}048-bit Morgan fingerprints \cite{morgan1965,rogers2010}. Median maxima across the
five folds run from 0.485 to 0.541 under the group split and from 0.630 to 0.644 under the
matched random split for Docking-44, and from 0.467 to 0.472 against 0.491 to 0.493 for
DOCKSTRING-58
(\path{results/public_scaffold_holdout_recovery/heldout_maximum_morgan_tanimoto_summary.csv}).
Close analogues survive the split: in the worst Docking-44 fold, 3.0\% of held-out ligands still
have a calibration neighbour at Tanimoto 1, which is bit-vector identity on this fingerprint and
not molecular identity. The split is not fingerprint-cluster-disjoint, and it measures recovery
inside one source collection rather than transfer to an external library.

\sisubsection{S5.3 Exploratory five-mode compression}

This analysis runs on the 20 DOCKSTRING kinase columns shared with the assay panels, named in
\texttt{provenance/targets} of \path{results/biological_core_modes/summary.json}, over the 190
target pairs they define. Of the 260{,}060 complete DOCKSTRING rows we removed the 481 whose
14-character connectivity block appears in an assay panel, leaving 259{,}579, and clipped 22
positive scores to zero. A conservative variant also removes the 3{,}039 rows sharing a cyclic
Bemis--Murcko scaffold with an assay compound, leaving 256{,}540. KiRHub supplies compound names
without structures, so neither exclusion can be applied against it.

The endpoint here is not the frozen endpoint of S5.1. Each panel supplies its own positives: the
upper decile of that panel's centred experimental target-pair correlations, 19 of 190 for every
panel. Rank agreement is the Spearman correlation between the docking and the experimental edge
vectors, and the retrieval metrics score the docking edges against those 19 labels.

\begin{table}[htbp]
\centering
\caption{The retained residual eigenvalues on the 20-kinase DOCKSTRING support. The correlation
trace is 20 and 19 eigenvalues are positive. Sources:
\texttt{results/biological\_core\_modes/retained\_eigenvalue\_spectrum.csv} for the eigenvalues,
and \texttt{selection\_rule/full\_eigenvalues} in the same directory's \texttt{summary.json} for
the cumulative fractions.}
\label{tab:s5spectrum}
\small
\begin{tabular}{rrr}
\toprule
Mode & Eigenvalue & Cumulative fraction of trace \\
\midrule
1 & 4.130 & 0.207 \\
2 & 2.133 & 0.313 \\
3 & 1.732 & 0.400 \\
4 & 1.201 & 0.460 \\
5 & 1.140 & 0.517 \\
\bottomrule
\end{tabular}
\end{table}

The selection rule takes the smallest $k$ whose leading residual eigenvalues reach half the
residual trace. Table~\ref{tab:s5spectrum} shows why that $k$ is five: four modes reach 0.460 of
the trace and five reach 0.517. The rule reads the docking eigenvalues alone. Under the
conservative scaffold exclusion it returns the same $k$, at cumulative fraction 0.517. The
participation ratio of the same 20-target residual map is 11.62
(\texttt{local\_20\_correlation\_PR} in \path{results/manuscript_evidence.json}), which is 2.3
times the number of retained modes.

Stability under a scaffold-disjoint split is in
\path{results/biological_core_modes/scaffold_support_stability.csv}. Whole Bemis--Murcko groups
go to five folds by SHA-256 assignment, with acyclic molecules grouped by recomputed InChIKey
connectivity block. Folds hold 50{,}256 to 55{,}261 ligands over 19{,}800 to 20{,}323 groups. The
rule returns $k=5$ in every fold, at cumulative fractions from 0.514 to 0.520. Measured against
the all-rows subspace, a fold's five-mode subspace has mean squared canonical correlation from
0.996 to 0.998 and maximum principal angle from 3.6 to 7.7 degrees. Across the ten fold pairs
(\path{results/biological_core_modes/scaffold_support_pairwise_stability.csv}) the same
quantities run from 0.986 to 0.996 and from 6.5 to 13.8 degrees.

\begin{table}[htbp]
\centering
\caption{Agreement of the docking target map with each experimental panel, before and after
truncating the residual correlation to five modes and renormalising to unit diagonal. Source: the
\texttt{exact\_connectivity\_excluded} rows of
\texttt{results/biological\_core\_modes/panel\_metrics.csv}. PKIS1 was the discovery panel.}
\label{tab:s5panels}
\small
\begin{tabular}{lrrrrrrrrr}
\toprule
& \multicolumn{3}{c}{Rank agreement} & \multicolumn{3}{c}{ROC-AUC} & \multicolumn{3}{c}{Average precision} \\
\cmidrule(lr){2-4}\cmidrule(lr){5-7}\cmidrule(lr){8-10}
Panel & Full & 5 modes & $\Delta$ & Full & 5 modes & $\Delta$ & Full & 5 modes & $\Delta$ \\
\midrule
PKIS1 & 0.174 & 0.210 & $+0.036$ & 0.664 & 0.677 & $+0.014$ & 0.300 & 0.265 & $-0.035$ \\
DAVIS & 0.280 & 0.326 & $+0.046$ & 0.755 & 0.791 & $+0.036$ & 0.340 & 0.375 & $+0.034$ \\
PKIS2 & 0.287 & 0.320 & $+0.033$ & 0.713 & 0.730 & $+0.017$ & 0.295 & 0.268 & $-0.027$ \\
KiRHub & 0.322 & 0.379 & $+0.058$ & 0.744 & 0.807 & $+0.063$ & 0.289 & 0.333 & $+0.044$ \\
\bottomrule
\end{tabular}
\end{table}

Table~\ref{tab:s5panels} gives the per-panel deltas. Rank agreement and ROC-AUC rise on all four
panels. Average precision splits: it rises on DAVIS and KiRHub and falls on PKIS1 and PKIS2.
Under the conservative scaffold exclusion the same pattern holds, with every delta within 0.010
of the value shown. The largest movement is the DAVIS average-precision delta, which goes from
$+0.034$ to $+0.025$.

\begin{table}[htbp]
\centering
\caption{Single-target deletions. Each panel is re-evaluated after dropping one of the 20
kinases, and the entry counts how many of the 20 deletions give a positive
five-mode-minus-full delta, with the median delta in parentheses. Source:
\texttt{target\_jackknife} in
\texttt{results/biological\_core\_modes/summary.json} and the same directory's
\texttt{target\_jackknife.csv}.}
\label{tab:s5jack}
\small
\begin{tabular}{lccc}
\toprule
Panel & Rank agreement & ROC-AUC & Average precision \\
\midrule
PKIS1 & 18/20 ($+0.037$) & 18/20 ($+0.018$) & 8/20 ($-0.002$) \\
DAVIS & 19/20 ($+0.043$) & 20/20 ($+0.037$) & 16/20 ($+0.043$) \\
PKIS2 & 19/20 ($+0.034$) & 20/20 ($+0.019$) & 5/20 ($-0.014$) \\
KiRHub & 19/20 ($+0.055$) & 18/20 ($+0.052$) & 18/20 ($+0.047$) \\
\bottomrule
\end{tabular}
\end{table}

Table~\ref{tab:s5jack} makes the split explicit. The average-precision gain is positive in only 5
of 20 deletions on PKIS2 and 8 of 20 on PKIS1, and its median delta is negative on both panels.
The selection rule survives every deletion: $k=5$ in all 80 panel-by-deletion rows, at cumulative
fractions from 0.508 to 0.541.

On the three panels scored without retuning, a quadratic assignment test that applies 50{,}000
complete relabellings of the 20 docking targets gives a mean five-mode-minus-full delta of 0.045
for rank agreement at one-sided probability 0.013, and 0.039 for ROC-AUC at probability 0.013.
The average-precision mean delta is 0.017 at probability 0.080. Holm correction across the three
metrics gives 0.038, 0.038 and 0.080; correcting across both low-mode representations, six tests
in total, gives 0.063, 0.063 and 0.159
(\path{results/biological_core_modes/omnibus_validation_qap.csv}).

The trace rule read no experimental value, and the three evaluation panels were scored without
retuning it. That does not make the comparison a clean test. The 50\% convention was fixed after
these resources had been inspected, and PKIS1 served as the discovery panel. We treat the result
as post hoc exploratory evidence that needs a new external panel.
